% ARCHIVED at v3.95 (2026-09-24) -- the 6.4 tables 6.17-6.19 (text-heavy) replaced by key points (author).
% Verbatim, not part of the build.

%% ---- 6.4.1-6.4.3 with Tables 6.17-6.19 (v3.93) ----
\subsection{Generative Models}\label{sec:res:summary:models}

The results fall into the lines the study was built on (\autoref{tab:summary-models}). D3IL-avoiding carries
the central result: the average-velocity models Pareto-dominate the diffusion model of \ac{DPCC} on its own
benchmark and with its own temporal U-Net. D3IL-aligning repeats it on camera observations and an alignment
through contact. On the quadrotor, UAV-pillars reaffirms D3IL-avoiding when its plans are flown and shows
what the change of plant costs, and UAV-corridor shows that the concept of \ac{DPCC} works in a new,
three-dimensional environment and that on straight-line demonstrations the generative model does not matter.
UAV-s-curve takes the plans onto a route of sharp turns, where the controller that flies them matters as much as
the model that makes them.

\begin{table}[htpb]
  \caption[Generative models across environments]{The conclusion of each environment on the generative models,
  as found in Sections~\ref{sec:res:state}--\ref{sec:res:uav}. Temporal U-Net of 3.96--4.0\,M parameters
  throughout; D3IL-aligning on training-set contexts.}
  \label{tab:summary-models}
  \centering
  \small
  \begin{tabularx}{\linewidth}{@{}p{2.2cm} L@{}}
    \toprule
    Environment & Conclusion \\
    \midrule
    D3IL-avoiding & the central result, Pareto dominance over the baseline: CI-MeanFM at $\nfe=1$ holds S\&C
                    $1.000$ in $59.2$ control steps at $18.1$\,ms, against the baseline's $1.000$ in $70.1$ at
                    $553.4$\,ms; MeanFM trades one episode in thirty ($0.967$, $58.6$ steps, $18.3$\,ms); FM at
                    $\nfe=1$ dominates the baseline too ($1.000$, $67.0$ steps, $17.3$\,ms); against the baseline
                    retrained at two steps ($1.000$, $61.2$ steps, $211.4$\,ms) the saving is twelvefold \\
    D3IL-aligning & the result repeated on camera observations and contact, every non-dominated configuration
                    MeanFM: under per-step projection it
                    dominates the baseline's best projected row, nine violation-free contexts against four,
                    $0.179$ against $0.462$\,m, $266$ against $809$\,ms; unprojected at $\nfe=20$ it closes
                    $84\,\%$ of the distance, against $14$, $10$ and $-2\,\%$ for CI-MeanFM, FM and the baseline \\
    UAV-pillars   & D3IL-avoiding reaffirmed on the quadrotor: the declared constraints transfer exactly and the
                    two average-velocity models keep their order; the vehicle tracks more closely than the arm,
                    and collisions with pillars outside the constraint sets end a third of the flights (S\&C
                    $0.600$--$0.700$) \\
    UAV-corridor  & new: the concept of \ac{DPCC} works in three dimensions, the violating steps of a flight
                    falling from $29$--$99$ to $0$--$10$; on straight-line demonstrations the three objectives
                    cannot be told apart and the budget decides; against the baseline a trade-off, $1/24$ and
                    $1/9$ of its time for more violating steps \\
    UAV-s-curve   & the controller matters: the cascaded geometric controller succeeds as often as MuJoCo MPC
                    without projection and on $5/10$ against $0/10$ with it, at a twenty-fourth of the cost; FM
                    at $\nfe=1$ leads on this single, simple route, and more evaluations are not always better
                    where a flight is easy to lose \\
    \bottomrule
  \end{tabularx}
\end{table}

\subsection{Projection Methods}\label{sec:res:constraints}

Endpoint projection beats the per-step projection of \ac{DPCC} where the task is run at a budget that gives
it steps to guide, and has nothing to win where one or two evaluations solve the task
(\autoref{tab:summary-projection}). On D3IL-aligning, at twenty evaluations, it keeps all ten contexts free
of violations. On UAV-corridor it overtakes the per-step projector over the hump at twenty evaluations with a
single candidate, and under the tilt it leaves fewer violating steps only at more time. On D3IL-avoiding it
is level or ahead at three and ten evaluations, but the task is solved at one, where it has no guiding step,
and UAV-pillars flies that operating point under per-step projection. On UAV-s-curve, at one evaluation, only
per-step projection runs, and it costs crossings without removing the violations the plans carry.

\begin{table}[htpb]
  \caption[Projection methods across environments]{Endpoint against per-step projection per environment, as
  found in Sections~\ref{sec:res:state}--\ref{sec:res:uav}. \emph{Budget}: the network evaluations at which
  the environment is operated or its projectors compared.}
  \label{tab:summary-projection}
  \centering
  \small
  \begin{tabularx}{\linewidth}{@{}p{2.2cm} r L@{}}
    \toprule
    Environment & Budget & Endpoint against per-step projection \\
    \midrule
    D3IL-avoiding & $1$      & no guiding step at the operating budget; at $3$ and $10$, with matched candidates,
                               level or ahead on the constraint and cheaper on every model except FM at $3$ \\
    D3IL-aligning & $20$     & ahead: all ten contexts free of violations and the box moved in eight, at the time
                               per control step of per-step projection; more contexts violation-free in seven of
                               nine comparisons \\
    UAV-pillars   & $1$, $2$ & per-step projection only, at the D3IL-avoiding operating points: as there, it turns
                               plans that violate the declared constraints into plans that violate none \\
    UAV-corridor  & $1$--$20$ & ahead only over the hump at $20$ with one candidate, $0.2$ against $1.8$ violating
                               steps at $291$ against $355$\,ms; under the tilt at $20$ fewer violating steps at
                               more time; at $3$ and $5$ as many or more \\
    UAV-s-curve   & $1$      & per-step projection only, endpoint projection having no guiding step: it binds the
                               measured position, leaves the corner cut of the plans in place and costs four of the
                               nine crossings, at $126$--$155$\,ms of projection per step \\
    \bottomrule
  \end{tabularx}
\end{table}

\subsection{Conclusion}\label{sec:res:summary:conclusion}

Read across the environments, the thesis rests on the two D3IL tasks. On D3IL-avoiding the average-velocity
models Pareto-dominate the diffusion model of \ac{DPCC}, with its own temporal U-Net, in fewer control steps
at a thirtieth of its time per action; on D3IL-aligning analytic average-velocity matching does so again,
closer to the target at a third of the time, and endpoint projection is the better projector on the
constraint there. The
quadrotor scenes carry the result to a new plant. UAV-pillars reaffirms D3IL-avoiding when its plans are
flown, and shows that a collision with an obstacle the constraints leave out ends a quadrotor's flight.
UAV-corridor shows that the concept of \ac{DPCC} works in a three-dimensional environment against constraints
no demonstration satisfies, and that on demonstrations as simple as straight lines the generative model does
not matter and the budget does. For the second research question the budget decides: endpoint projection
beats the per-step projection of \ac{DPCC} at twenty evaluations, on D3IL-aligning and over the hump of
UAV-corridor, and has nothing to guide where one or two evaluations solve the task. UAV-s-curve adds the controller: in
this framework the same plans succeed or fail by the controller that flies them, the cascaded geometric
controller used throughout the thesis is the better of the two on the configuration of record, and constraint
satisfaction is the plan's, since every learned plan cuts a corner its demonstrations clear.

Three findings cut across the environments. The price of a projection is set by what the projector is
handed: a plan close to the feasible set is cheap to repair, which is why the flow-based models at one
evaluation and endpoint projection's predicted endpoint are cheap, and the two-step diffusion baseline,
handed a plan one step from noise, is not. A task saturated at one evaluation hides what a task that rewards
more evaluations exposes: the self-referential target of the consistency-interpolated model, invisible on
D3IL-avoiding, stops it improving on D3IL-aligning. And the demonstrations bound what a model can show: on
the straight lanes of UAV-corridor the objectives do not differ, and the plans stay within the altitudes the
demonstrations cover, which is where per-step projection at the smallest budgets falls short.
\autoref{tab:summary-combinations} collects the combination each environment supports.

\begin{table}[htpb]
  \caption[Selected model--projection combinations]{The supported combination per environment.}
  \label{tab:summary-combinations}
  \centering
  \small
  \begin{tabularx}{\linewidth}{@{}p{2.2cm} L L@{}}
    \toprule
    Environment & Combination & Result \\
    \midrule
    D3IL-avoiding
      & MeanFM or CI-MeanFM, $\nfe=1$, per-step projection, temporal consistency
      & S\&C $0.967$ and $1.000$ in $58.6$ and $59.2$ control steps at $18.3$ and $18.1$\,ms per step, against the
        baseline's $1.000$ in $70.1$ at $553.4$\,ms; five seeds, the protocol of \ac{DPCC} \\
    D3IL-aligning
      & MeanFM, $\nfe=20$, endpoint projection, random selection
      & $10/10$ contexts free of violations, box moved in $8/10$ \\
    UAV-pillars
      & the D3IL-avoiding configurations: MeanFM or CI-MeanFM, $\nfe=1$ or $2$, per-step projection, temporal
        consistency
      & no violating step; S\&C $0.600$ to $0.700$, every failure a collision with a pillar outside the
        constraint sets \\
    UAV-corridor
      & any flow-based objective, the budget chosen for the trade
      & violating steps per flight from $10.2$ at $27.4$\,ms ($\nfe=1$, per-step) to $1.4$ at $416.9$\,ms
        ($\nfe=20$, endpoint) under the tilt, from $8.0$ at $71.9$\,ms ($\nfe=3$, endpoint) to $0.2$ at
        $291.0$\,ms ($\nfe=20$, endpoint, one candidate) over the hump; every flight reaches the end \\
    UAV-s-curve
      & FM, $\nfe=1$, without projection, the cascaded geometric controller
      & nine of ten flights cross, none free of violations; MuJoCo MPC crosses as often without projection and on
        none with it, at twenty-four times the controller's cost per step \\
    \bottomrule
  \end{tabularx}
\end{table}

